\documentclass[11pt,a4paper]{article}

% ── Packages ───────────────────────────────────────────────────
\usepackage[utf8]{inputenc}
\usepackage[T1]{fontenc}
\usepackage{lmodern}
\usepackage[margin=2.5cm]{geometry}
\usepackage{xcolor}
\usepackage{hyperref}
\usepackage{listings}
\usepackage{booktabs}
\usepackage{enumitem}
\usepackage{titlesec}
\usepackage{fancyhdr}
\usepackage{tcolorbox}
\usepackage{parskip}
\usepackage{longtable}

% ── Colours ────────────────────────────────────────────────────
\definecolor{fridaybrand}{HTML}{1E5A96}
\definecolor{codebg}{HTML}{F5F2EB}
\definecolor{codefg}{HTML}{23211E}
\definecolor{getgreen}{HTML}{2E7D32}
\definecolor{postorange}{HTML}{E65100}
\definecolor{sectionblue}{HTML}{1565C0}

% ── Hyperlinks ─────────────────────────────────────────────────
\hypersetup{
  colorlinks  = true,
  linkcolor   = fridaybrand,
  urlcolor    = fridaybrand,
}

% ── Code listings ──────────────────────────────────────────────
\lstset{
  backgroundcolor = \color{codebg},
  basicstyle      = \ttfamily\small\color{codefg},
  breaklines      = true,
  frame           = single,
  rulecolor       = \color{fridaybrand!30},
  tabsize         = 2,
  showstringspaces= false,
  captionpos      = b,
}
\lstdefinelanguage{json}{
  morestring=[b]",
  literate=
    *{0}{{{\color{blue}0}}}{1}
     {1}{{{\color{blue}1}}}{1}
     {2}{{{\color{blue}2}}}{1}
     {3}{{{\color{blue}3}}}{1}
     {4}{{{\color{blue}4}}}{1}
     {5}{{{\color{blue}5}}}{1}
     {6}{{{\color{blue}6}}}{1}
     {7}{{{\color{blue}7}}}{1}
     {8}{{{\color{blue}8}}}{1}
     {9}{{{\color{blue}9}}}{1}
}

% ── Section styling ────────────────────────────────────────────
\titleformat{\section}
  {\Large\bfseries\color{sectionblue}}{\thesection}{1em}{}
\titleformat{\subsection}
  {\large\bfseries\color{fridaybrand}}{\thesubsection}{1em}{}
\titleformat{\subsubsection}
  {\normalsize\bfseries\color{fridaybrand!80}}{\thesubsubsection}{1em}{}

% ── Header / Footer ───────────────────────────────────────────
\pagestyle{fancy}
\fancyhf{}
\fancyhead[L]{\textcolor{fridaybrand}{\textbf{Project Friday}}}
\fancyhead[R]{\textcolor{gray}{API Reference}}
\fancyfoot[C]{\thepage}
\renewcommand{\headrulewidth}{0.4pt}

% ── Callout boxes ──────────────────────────────────────────────
\newtcolorbox{notebox}{
  colback=blue!5, colframe=fridaybrand, title=Note,
  fonttitle=\bfseries, arc=2mm
}
\newtcolorbox{warningbox}{
  colback=red!5, colframe=red!60!black, title=Warning,
  fonttitle=\bfseries, arc=2mm
}

% ── HTTP method badges ─────────────────────────────────────────
\newcommand{\httpget}{\textcolor{getgreen}{\textbf{GET}}}
\newcommand{\httppost}{\textcolor{postorange}{\textbf{POST}}}

% ══════════════════════════════════════════════════════════════
\begin{document}

\begin{titlepage}
\centering
\vspace*{3cm}
{\Huge\bfseries\color{fridaybrand} Project Friday\par}
\vspace{0.6cm}
{\Large API Reference\par}
\vspace{1.5cm}
{\large Backend REST endpoints, SSE event protocol,\\tool specifications, and configuration reference.\par}
\vspace{2cm}
{\normalsize Version 0.2.0\par}
{\normalsize \today\par}
\vfill
\end{titlepage}

\tableofcontents
\newpage

% ══════════════════════════════════════════════════════════════
\section{Backend REST API}

Base URL: \texttt{http://localhost:8000}

% ── POST /chat ─────────────────────────────────────────────────
\subsection{\httppost{} /chat}

Sends a user query to the Friday agent and streams the response as Server-Sent Events (SSE).

\subsubsection{Request}

\begin{lstlisting}[language=json,title={Request body}]
{
  "query": "Create a Next.js project called Gargantua",
  "conversation_id": "optional-uuid"
}
\end{lstlisting}

\begin{table}[h!]
\centering
\begin{tabular}{@{}llll@{}}
\toprule
\textbf{Field} & \textbf{Type} & \textbf{Required} & \textbf{Description} \\
\midrule
\texttt{query} & string & Yes & The user's message \\
\texttt{conversation\_id} & string & No & UUID to continue a conversation. Auto-generated if omitted. \\
\bottomrule
\end{tabular}
\end{table}

\subsubsection{Response}

Content-Type: \texttt{text/event-stream}

Each SSE frame has the format:

\begin{lstlisting}[language=json,title={SSE data frame}]
data: {"type": "thought|tool_call|tool_result|final", "content": "..."}
\end{lstlisting}

\begin{table}[h!]
\centering
\begin{tabular}{@{}lp{9cm}@{}}
\toprule
\textbf{Event Type} & \textbf{Description} \\
\midrule
\texttt{thought} & Internal reasoning step from the LLM (before tool calls) \\
\texttt{tool\_call} & Tool being invoked: ``\texttt{tool\_name \{args\_json\}}'' \\
\texttt{tool\_result} & Tool output: ``\texttt{[tool\_name] output text}'' \\
\texttt{final} & Agent's final answer to the user \\
\bottomrule
\end{tabular}
\end{table}

\subsubsection{Error Handling}

If the agent encounters a fatal error, a final SSE frame is sent:

\begin{lstlisting}[language=json]
data: {"type": "final", "content": "Friday failed: <error message>"}
\end{lstlisting}

% ── GET /workspace ─────────────────────────────────────────────
\subsection{\httpget{} /workspace}

Returns a recursive file tree of the sandboxed workspace directory.

\subsubsection{Response}

\begin{lstlisting}[language=json,title={Response body}]
{
  "tree": [
    {
      "name": "src",
      "path": "src",
      "type": "directory",
      "children": [
        {"name": "index.ts", "path": "src/index.ts", "type": "file"}
      ]
    },
    {"name": "package.json", "path": "package.json", "type": "file"}
  ]
}
\end{lstlisting}

\begin{table}[h!]
\centering
\begin{tabular}{@{}llp{7cm}@{}}
\toprule
\textbf{Field} & \textbf{Type} & \textbf{Description} \\
\midrule
\texttt{name} & string & File or directory name \\
\texttt{path} & string & Workspace-relative path (forward slashes) \\
\texttt{type} & string & \texttt{"file"} or \texttt{"directory"} \\
\texttt{children} & array & Nested entries (directories only) \\
\bottomrule
\end{tabular}
\end{table}

% ── GET /skills ────────────────────────────────────────────────
\subsection{\httpget{} /skills}

Returns all committed skill scripts from the library.

\begin{lstlisting}[language=json,title={Response body}]
{
  "skills": [
    {
      "name": "hn_scraper",
      "description": "Scrapes Hacker News front page",
      "path": "skills/hn_scraper.py"
    }
  ]
}
\end{lstlisting}

% ── GET /agents ────────────────────────────────────────────────
\subsection{\httpget{} /agents}

Returns all registered Skill Agents.

\begin{lstlisting}[language=json,title={Response body}]
{
  "agents": [
    {
      "name": "nextjs",
      "context_preview": "## Skill Agent: Next.js\n**Description:** ..."
    }
  ]
}
\end{lstlisting}

% ══════════════════════════════════════════════════════════════
\section{Frontend Proxy Routes}

The Next.js frontend proxies all backend requests through \texttt{/api/friday/*} to avoid CORS issues.

\begin{table}[h!]
\centering
\begin{tabular}{@{}llll@{}}
\toprule
\textbf{Method} & \textbf{Frontend Path} & \textbf{Backend Path} & \textbf{Description} \\
\midrule
POST & \texttt{/api/friday} & \texttt{/chat} & Chat SSE stream \\
GET & \texttt{/api/friday} & \texttt{/workspace} & File tree \\
GET & \texttt{/api/friday/skills} & \texttt{/skills} & Skill scripts \\
GET & \texttt{/api/friday/agents} & \texttt{/agents} & Skill Agents \\
\bottomrule
\end{tabular}
\end{table}

% ══════════════════════════════════════════════════════════════
\section{Tool Specifications}

\subsection{File Operations}

\subsubsection{list\_files}

\begin{tabular}{@{}lp{10cm}@{}}
\textbf{Arg} & \texttt{directory: str = "."} --- workspace-relative path \\
\textbf{Returns} & Newline-separated list of entries, or \texttt{(empty directory)} \\
\end{tabular}

\subsubsection{read\_file}

\begin{tabular}{@{}lp{10cm}@{}}
\textbf{Arg} & \texttt{filename: str} --- workspace-relative file path \\
\textbf{Returns} & Full file contents as string \\
\end{tabular}

\subsubsection{write\_to\_file}

\begin{tabular}{@{}lp{10cm}@{}}
\textbf{Args} & \texttt{filename: str}, \texttt{content: str} \\
\textbf{Returns} & ``Wrote N chars to filename'' \\
\textbf{Note} & Creates parent directories automatically \\
\end{tabular}

\subsubsection{append\_to\_file}

\begin{tabular}{@{}lp{10cm}@{}}
\textbf{Args} & \texttt{filename: str}, \texttt{content: str} \\
\textbf{Returns} & ``Appended N chars to filename'' \\
\end{tabular}

\subsubsection{create\_directory}

\begin{tabular}{@{}lp{10cm}@{}}
\textbf{Arg} & \texttt{directory: str} --- workspace-relative path \\
\textbf{Returns} & ``Directory created: dirname'' \\
\end{tabular}

\subsubsection{delete\_file}

\begin{tabular}{@{}lp{10cm}@{}}
\textbf{Arg} & \texttt{filename: str} --- workspace-relative file path \\
\textbf{Returns} & ``Deleted: filename'' or error message \\
\textbf{Note} & Only deletes files, not directories \\
\end{tabular}

\subsection{Shell Execution}

\subsubsection{execute\_bash\_command}

\begin{tabular}{@{}lp{10cm}@{}}
\textbf{Arg} & \texttt{command: str} --- shell command to run \\
\textbf{CWD} & \texttt{WORKSPACE\_DIR} \\
\textbf{Timeout} & \texttt{COMMAND\_TIMEOUT} (default 120s) \\
\textbf{Validation} & First token must be on allowlist; no chaining operators \\
\textbf{Returns} & stdout on success, \texttt{ERROR: ...} on failure or block \\
\end{tabular}

\subsubsection{execute\_in\_directory}

\begin{tabular}{@{}lp{10cm}@{}}
\textbf{Args} & \texttt{directory: str} (absolute path), \texttt{command: str} \\
\textbf{Validation} & Directory must be under \texttt{ALLOWED\_TARGET\_DIRS}; command must be on allowlist \\
\textbf{Returns} & Same as \texttt{execute\_bash\_command} \\
\end{tabular}

\subsection{Web Search}

\subsubsection{web\_search}

\begin{tabular}{@{}lp{10cm}@{}}
\textbf{Arg} & \texttt{query: str} --- search query \\
\textbf{Provider} & Tavily (if \texttt{TAVILY\_API\_KEY} set), else DuckDuckGo \\
\textbf{Returns} & Top 3 result summaries joined by blank lines \\
\end{tabular}

\subsection{Skill Library}

\subsubsection{save\_to\_skill\_library}

\begin{tabular}{@{}lp{10cm}@{}}
\textbf{Args} & \texttt{skill\_name: str}, \texttt{description: str}, \texttt{temp\_filename: str} \\
\textbf{Action} & Moves workspace file to \texttt{skills/<slug>.py}, updates \texttt{index.json} \\
\textbf{Returns} & ``Skill 'name' committed to library.'' \\
\end{tabular}

\subsubsection{search\_skill\_library}

\begin{tabular}{@{}lp{10cm}@{}}
\textbf{Arg} & \texttt{query: str} --- description to match against \\
\textbf{Returns} & JSON with \texttt{name}, \texttt{description}, \texttt{path}, \texttt{similarity} \\
\end{tabular}

\subsection{Skill Agents}

\subsubsection{create\_skill\_agent}

\begin{tabular}{@{}lp{10cm}@{}}
\textbf{Args} & \texttt{framework\_name}, \texttt{description}, \texttt{dos} (newline-separated), \texttt{donts} (newline-separated), \texttt{scaffold\_steps} (newline-separated), \texttt{trigger\_patterns} (comma-separated) \\
\textbf{Action} & Creates \texttt{skills/agents/<slug>/} with \texttt{manifest.json}, \texttt{style\_guide.md}, \texttt{templates/} \\
\textbf{Returns} & Summary of created agent \\
\end{tabular}

\subsubsection{load\_skill\_context}

\begin{tabular}{@{}lp{10cm}@{}}
\textbf{Arg} & \texttt{skill\_name: str} \\
\textbf{Returns} & Text block with dos, don'ts, scaffold steps, and full style guide \\
\end{tabular}

\subsubsection{list\_skill\_agents}

\begin{tabular}{@{}lp{10cm}@{}}
\textbf{Args} & None \\
\textbf{Returns} & Formatted list of all registered agents \\
\end{tabular}

\subsubsection{update\_skill\_agent}

\begin{tabular}{@{}lp{10cm}@{}}
\textbf{Args} & \texttt{skill\_name: str}, \texttt{updates: str} (JSON) \\
\textbf{JSON keys} & \texttt{add\_dos}, \texttt{add\_donts}, \texttt{remove\_dos}, \texttt{remove\_donts}, \texttt{add\_scaffold\_steps}, \texttt{description} \\
\textbf{Returns} & Summary of updated counts \\
\end{tabular}

% ══════════════════════════════════════════════════════════════
\section{Configuration Reference}

\begin{longtable}{@{}llp{6cm}@{}}
\toprule
\textbf{Variable} & \textbf{Default} & \textbf{Description} \\
\midrule
\endhead
\texttt{MODEL\_PROVIDER} & \texttt{lmstudio} & LLM provider: \texttt{lmstudio}, \texttt{groq}, \texttt{ollama} \\
\texttt{LMSTUDIO\_BASE\_URL} & \texttt{http://localhost:1234/v1} & LM Studio API endpoint \\
\texttt{LMSTUDIO\_MODEL} & \texttt{local-model} & Requested model (auto-detected if \texttt{LMSTUDIO\_AUTO\_SELECT=true}) \\
\texttt{LMSTUDIO\_API\_KEY} & \texttt{lm-studio} & API key for LM Studio \\
\texttt{LMSTUDIO\_TEMPERATURE} & \texttt{0.7} & Sampling temperature \\
\texttt{LMSTUDIO\_AUTO\_SELECT} & \texttt{true} & Auto-detect loaded model \\
\texttt{GROQ\_API\_KEY} & (empty) & Groq cloud API key \\
\texttt{GROQ\_MODEL} & \texttt{llama-3.3-70b-versatile} & Groq model name \\
\texttt{OLLAMA\_MODEL} & \texttt{qwen2.5-coder:7b} & Ollama model name \\
\texttt{MAX\_TOKENS} & \texttt{2048} & Max tokens per LLM response \\
\texttt{TAVILY\_API\_KEY} & (empty) & Tavily web search API key \\
\texttt{WORKSPACE\_DIR} & \texttt{./workspace} & Sandboxed file-system root \\
\texttt{SKILLS\_DIR} & \texttt{./skills} & Skill library directory \\
\texttt{MAX\_TOOL\_ATTEMPTS} & \texttt{3} & Retry budget for failed tools \\
\texttt{RECURSION\_LIMIT} & \texttt{25} & LangGraph recursion hard cap \\
\texttt{COMMAND\_TIMEOUT} & \texttt{120} & Max subprocess runtime (seconds) \\
\texttt{COMMAND\_ALLOWLIST} & \texttt{python,npm,...} & Allowed command prefixes (comma-separated) \\
\texttt{ALLOWED\_TARGET\_DIRS} & (empty) & Directories for external scaffolding (comma-separated absolute paths) \\
\texttt{NEXTJS\_ORIGIN} & \texttt{http://localhost:3000} & Frontend origin for CORS \\
\texttt{FRIDAY\_BACKEND\_URL} & \texttt{http://localhost:8000} & Backend URL (used by frontend proxy) \\
\bottomrule
\end{longtable}

% ══════════════════════════════════════════════════════════════
\section{Error Codes \& Responses}

\begin{table}[h!]
\centering
\begin{tabular}{@{}llp{7cm}@{}}
\toprule
\textbf{Scenario} & \textbf{Response} & \textbf{Details} \\
\midrule
Path escape attempt & \texttt{PermissionError} & ``Path escape attempt blocked.'' \\
Absolute path in workspace & \texttt{PermissionError} & ``Absolute paths are not allowed.'' \\
Symlink to external path & \texttt{PermissionError} & ``Symbolic link target is outside the workspace.'' \\
Blocked path segment & \texttt{PermissionError} & ``Access to protected path segment blocked.'' \\
Command not on allowlist & \texttt{BLOCKED:} string & Lists the command and allowed alternatives \\
Shell chaining detected & \texttt{BLOCKED:} string & ``Shell operator 'X' is not allowed.'' \\
Command timeout & \texttt{ERROR:} string & ``Command timed out after Ns.'' \\
External dir not allowed & \texttt{BLOCKED:} string & Lists allowed target directories \\
Backend unreachable & HTTP 502 & Frontend proxy returns error JSON \\
\bottomrule
\end{tabular}
\end{table}

\end{document}
